\chapter{Hochschild Calculus for CY Categories}
\label{ch:hochschild-calculus}

A Calabi--Yau $d$-category has a non-degenerate Serre pairing. What does this pairing force on the Hochschild invariants? For a general dg category, Hochschild homology $\HH_\bullet(\cC)$ carries only a Connes $B$-operator and Hochschild cohomology $\HH^\bullet(\cC)$ carries only a Gerstenhaber bracket; the two are linked by a cap product, but no intrinsic pairing relates them. The CY structure supplies exactly this pairing, and the rest of the chapter is its consequence.

\begin{remark}[Three Hochschild theories]
\label{rem:hochschild-convention-categorical}
Three distinct Hochschild theories appear in the three volumes and must be kept separate.
\begin{enumerate}
 \item \emph{Categorical Hochschild} $\HH^\bullet_{\mathrm{cat}}$ (this chapter): input a dg category $\cC$; cyclic bar complex $\mathrm{CC}_\bullet(\cC)$; output $E_2$-algebra (Deligne); for $\cC$ a $\CY_d$ category, also a $(2{-}d)$-shifted Poisson / $\BV$ structure (Theorem~\textup{\ref{thm:shifted-poisson}}).
 \item \emph{Topological Hochschild} $\HH^\bullet_{\mathrm{top}}$: input an $E_1$-algebra over a base ring; output $E_2$-algebra (Deligne; Lurie). The Hochschild theory of Vol~II's slab algebra of 3d HT theories.
 \item \emph{Chiral Hochschild} $\ChirHoch^\bullet$: input an $E_\infty$-chiral algebra on a curve $X$; output concentrated in cohomological degrees $\{0,1,2\}$ (Vol~I Theorem~H). The chiral Gerstenhaber bracket on $\ChirHoch^\bullet$ is not the topological one.
\end{enumerate}
The CY-to-chiral correspondence programme $\{\Phi_d\}$ of Chapter~\textup{\ref{ch:cy-to-chiral}} sends the categorical input $\HH^\bullet_{\mathrm{cat}}(\cC)$ to a chiral output whose own chiral Hochschild cohomology $\ChirHoch^\bullet(\Phi_d(\cC))$ is sharper (concentrated in $\{0,1,2\}$). The categorical Gerstenhaber bracket on $\HH^\bullet_{\mathrm{cat}}$ maps to the chiral convolution bracket on $\ChirHoch^\bullet$ via the bar-cobar passage of Vol~I Chapter~\textup{\ref{chap:climax-platonic}}; this is not an identification of objects, only of structures. Throughout this chapter, $\HH^\bullet$ without qualifier means $\HH^\bullet_{\mathrm{cat}}$.
\end{remark}

\begin{remark}[$\kappa_\bullet$ convention in this chapter]
\label{rem:hochschild-kappa-convention}
Every $\kappa$ carries a subscript from $\{\mathrm{ch}, \mathrm{cat}, \mathrm{BKM}, \mathrm{fiber}\}$; the manifold invariants ($\kappa_{\mathrm{cat}}, \kappa_{\mathrm{fiber}}$) and algebraisation invariants ($\kappa_{\mathrm{ch}}, \kappa_{\mathrm{BKM}}$) split as in Chapter~\ref{ch:cy-categories}.
\end{remark}

The Serre pairing provides a quasi-isomorphism $\HH^\bullet(\cC) \simeq \HH_{d-\bullet}(\cC)$ of degree $-d$, so that Hochschild cohomology inherits the $B$-operator and Hochschild homology inherits the Gerstenhaber bracket. The combination is a $(2-d)$-shifted Poisson (equivalently, $\BV_{2-d}$) structure on $\HH^\bullet(\cC)$: under the PTVV convention an $n$-shifted Poisson bracket has cohomological degree $-n$, so the bracket $\{-,-\}_{2-d}$ has degree $d-2$, and the compatibility $[B, -] = \{-,-\}$ with $|B| = +1$ holds on the Hochschild complex. At $d = 2$ this is a $0$-shifted Poisson structure (ordinary Poisson, bracket of degree $0$); at $d = 3$, a $(-1)$-shifted Poisson structure (the CY$_3$ case of Pantev--To\"en--Vaqui\'e--Vezzosi, bracket of degree $+1$, i.e.\ a $\BV$ algebra); at $d = 1$, a $1$-shifted Poisson structure (bracket of degree $-1$) of the kind that drives factorization homology.

The shifted Poisson structure on $\HH^\bullet(\cC)$ is the categorical precursor of the convolution Lie bracket on the modular deformation complex of Volume~I: the Gerstenhaber bracket becomes the bar-cobar convolution bracket, the Connes $B$-operator becomes the modular differential, and the Serre pairing becomes the cyclic duality that controls the $\bS^d$-framing. The cyclic bar complex of Chapter~\ref{ch:cyclic-ainf} is the single-object specialization of the Hochschild complex constructed here.


\section{Hochschild homology and cohomology}
\label{sec:hh-basics}

The cyclic bar complex of a dg category is built from the multi-object generalisation of the cyclic $\Ainf$-complex of Chapter~\ref{ch:cyclic-ainf}, replacing a single algebra by tensor strings of morphisms $X_0 \to X_1 \to \cdots \to X_n \to X_0$. The Hochschild homology and cohomology defined below are the Morita-invariant invariants extracted from this complex.

\begin{definition}[Hochschild complex]
\label{def:hochschild-complex}
Let $\cC$ be a dg category over $k$. The \emph{Hochschild chain complex}
of $\cC$ is
\[
 \mathrm{CC}_n(\cC)
 = \bigoplus_{X_0, \ldots, X_n \in \mathrm{Ob}(\cC)}
 \Hom_\cC(X_{n}, X_0) \otimes
 \Hom_\cC(X_{n-1}, X_n) \otimes \cdots \otimes
 \Hom_\cC(X_0, X_1)
\]
with the Hochschild differential
\begin{align*}
 b(f_0 \otimes \cdots \otimes f_n)
 &= \sum_{i=0}^{n-1} (-1)^{\epsilon_i}\,
 f_0 \otimes \cdots \otimes f_i f_{i+1} \otimes \cdots \otimes f_n \\
 &\quad + (-1)^{\epsilon_n}\, f_n f_0 \otimes f_1 \otimes \cdots \otimes f_{n-1},
\end{align*}
where $\epsilon_i = |f_0| + \cdots + |f_i|$ (Koszul signs). The
\emph{Hochschild homology} is $\HH_\bullet(\cC) = H_\bullet(\mathrm{CC}_\bullet(\cC), b)$.
\end{definition}

\begin{definition}[Hochschild cohomology]
\label{def:hochschild-cohomology}
The \emph{Hochschild cochain complex} of a dg category $\cC$ is
\[
 \mathrm{CC}^n(\cC)
 = \prod_{X_0, \ldots, X_n \in \mathrm{Ob}(\cC)}
 \Hom_k\bigl(
 \Hom_\cC(X_{n-1}, X_n) \otimes \cdots \otimes \Hom_\cC(X_0, X_1),\;
 \Hom_\cC(X_0, X_n)
 \bigr)
\]
with the Hochschild coboundary $\delta$ dual to $b$. The
\emph{Hochschild cohomology} is
$\HH^\bullet(\cC) = H^\bullet(\mathrm{CC}^\bullet(\cC), \delta)$.
\end{definition}

\begin{remark}[Morita invariance]
\label{rem:hh-morita}
Both $\HH_\bullet$ and $\HH^\bullet$ are Morita invariant: if
$\cC$ and $\cD$ are derived Morita equivalent (i.e.,
$\Perf(\cC) \simeq \Perf(\cD)$ as dg categories), then
$\HH_\bullet(\cC) \simeq \HH_\bullet(\cD)$ and
$\HH^\bullet(\cC) \simeq \HH^\bullet(\cD)$. This is essential for
the CY programme, where different geometric models (Fukaya, derived
coherent sheaves, matrix factorizations) of the same CY category must
produce the same Hochschild invariants.
\end{remark}


\section{The Gerstenhaber bracket}
\label{sec:gerstenhaber-bracket}

The Hochschild cochain complex $\mathrm{CC}^\bullet(\cC)$ is a pre-Lie algebra under operadic composition of multilinear maps, and the pre-Lie commutator descends to a graded Lie bracket of degree $-1$ on $\HH^\bullet(\cC)$. Combined with the cup product, this is the Gerstenhaber algebra structure.

\begin{theorem}[Gerstenhaber structure on $\HH^\bullet$]
\label{thm:gerstenhaber}
\ClaimStatusProvedElsewhere
The Hochschild cohomology $\HH^\bullet(\cC)$ carries two compatible
structures:
\begin{enumerate}[label=(\roman*)]
 \item A graded commutative \emph{cup product}
 $\smile \colon \HH^p(\cC) \otimes \HH^q(\cC) \to \HH^{p+q}(\cC)$,
 the composition of multilinear maps;
 \item A degree $(-1)$ Lie bracket
 $[\cdot, \cdot] \colon
 \HH^p(\cC) \otimes \HH^q(\cC) \to \HH^{p+q-1}(\cC)$
 (the \emph{Gerstenhaber bracket}), defined at the cochain level by
 \[
 [f, g] = f \circ g - (-1)^{(|f|-1)(|g|-1)} g \circ f
 \]
 where $\circ$ denotes the operadic pre-Lie composition:
 \[
 (f \circ g)(a_1, \ldots, a_{p+q-1})
 = \sum_{i=1}^{p} (-1)^{\dagger_i}\,
 f(a_1, \ldots, a_{i-1}, g(a_i, \ldots, a_{i+q-1}), a_{i+q}, \ldots, a_{p+q-1}),
 \qquad \dagger_i = (i-1)(q-1) + (|g|-1)\sum_{j<i}(|a_j|-1),
 \]
\end{enumerate}
Together, the cup product and Gerstenhaber bracket make
$\HH^\bullet(\cC)$ a \emph{Gerstenhaber algebra}: the bracket is a
derivation of the cup product in each variable.
\end{theorem}

\begin{remark}[Gerstenhaber = $E_2$ cohomology]
\label{rem:gerstenhaber-e2}
By Kontsevich's formality theorem and the Deligne conjecture
(proved by Kontsevich--Soibelman, McClure--Smith, and others),
$\mathrm{CC}^\bullet(\cC)$ carries a natural $E_2$-algebra structure
whose cohomology-level shadow is the Gerstenhaber algebra
$(\HH^\bullet(\cC), \smile, [\cdot, \cdot])$. The $E_2$-structure
is the chain-level origin of the braided monoidal category
structure on $\Rep^{E_2}(\Zder(A))$ (Chapter~\ref{ch:drinfeld-center}).
\end{remark}

\begin{proposition}[Gerstenhaber bracket controls deformations]
\label{prop:gerstenhaber-deformations}
\ClaimStatusProvedElsewhere
The Gerstenhaber bracket controls the deformation theory of
$\cC$:
\begin{enumerate}[label=(\roman*)]
 \item $\HH^2(\cC)$ classifies first-order deformations of
 $\cC$ as a dg category;
 \item The obstruction to extending a first-order deformation
 $\mu_1 \in \HH^2(\cC)$ to second order lies in $\HH^3(\cC)$
 and is given by the Gerstenhaber square
 $\frac{1}{2}[\mu_1, \mu_1]$;
 \item A formal deformation is an element
 $\mu = \sum_{n \geq 1} \hbar^n \mu_n \in \HH^2(\cC)[[\hbar]]$
 satisfying the Maurer--Cartan equation
 $\delta \mu + \frac{1}{2}[\mu, \mu] = 0$.
\end{enumerate}
The MC equation for deformations of $\cC$ maps, under $\Phi$, to
the MC equation $D\Theta_A + \frac{1}{2}[\Theta_A, \Theta_A] = 0$
of the modular convolution algebra $\mathfrak{g}^{\mathrm{mod}}_A$ of Vol~I (the bar-cobar Lie algebra whose Maurer--Cartan elements parametrise modular characteristic corrections; Vol~I, thm:mc2-bar-intrinsic).
\end{proposition}


\section{The Connes $B$-operator and cyclic homology}
\label{sec:connes-b}

The Gerstenhaber structure is one half of the Hochschild calculus; the other half lives on homology and encodes the $S^1$-action on the cyclic bar complex. The Connes $B$-operator is the chain-level incarnation of this action.

\begin{definition}[Connes $B$-operator]
\label{def:connes-b}
The \emph{Connes operator}
$B \colon \mathrm{CC}_n(\cC) \to \mathrm{CC}_{n+1}(\cC)$ is the composite $B = (1 \otimes -) \circ N \circ (1 - t)$, where $t$ is the signed cyclic rotation and $N = \sum_{i=0}^n t^i$ the norm. Explicitly:
\[
 B(f_0 \otimes \cdots \otimes f_n)
 = \sum_{i=0}^{n} (-1)^{ni}\,
 1 \otimes f_i \otimes f_{i+1} \otimes \cdots
 \otimes f_{i-1}
\]
(indices mod $n+1$, with $1 \in \Hom_\cC(X_i, X_i)$ the identity endomorphism on the source object, inserted in the leading slot). It satisfies $B^2 = 0$ and the mixed-complex identity $bB + Bb = 0$.
\end{definition}

\begin{definition}[Cyclic, negative cyclic, and periodic cyclic homology]
\label{def:cyclic-homology-variants}
From the mixed complex $(\mathrm{CC}_\bullet(\cC), b, B)$
with $|b| = -1$, $|B| = +1$, $b^2 = B^2 = bB + Bb = 0$:
\begin{enumerate}[label=(\roman*)]
 \item \emph{Cyclic homology}:
 $\HC_\bullet(\cC) = H_\bullet(\mathrm{CC}_\bullet(\cC)[[u]] / u\mathrm{CC}_\bullet(\cC)[[u]], b + uB)$, i.e.\ the quotient by $u$ of the negative-cyclic complex (equivalently, the Connes bicomplex modulo the $1-t$ relation);
 \item \emph{Negative cyclic homology}:
 $\HC^-_\bullet(\cC) = H_\bullet(\mathrm{CC}_\bullet(\cC)[[u]], b + uB)$,
 where $|u| = 2$;
 \item \emph{Periodic cyclic homology}:
 $\HP_\bullet(\cC) = H_\bullet(\mathrm{CC}_\bullet(\cC)((u)), b + uB)$.
\end{enumerate}
The SBI exact triangle (Loday convention)
\[
 \cdots \longrightarrow \HH_n \xrightarrow{\; I \;} \HC_n \xrightarrow{\; S \;} \HC_{n-2} \xrightarrow{\; B \;} \HH_{n-1} \longrightarrow \cdots,
 \qquad |I|=0,\ |S|=-2,\ |B|=+1,
\]
relates the three theories: $I$ is induced by $\HC \otimes_{k[u]} k \to \HH$, $S$ is periodicity $\otimes u$, and $B$ is the Connes operator after passage to the quotient. The cumulative index shift around the triangle is $0 + (-2) + 1 = -1$, matching the periodic structure.
\end{definition}

\begin{example}[Smooth variety]
\label{ex:hkr}
For $\cC = \Perf(X)$ with $X$ a smooth algebraic variety of
dimension $n$, the Hochschild--Kostant--Rosenberg theorem gives
\[
 \HH_k(\Perf(X)) \simeq \bigoplus_{p-q=k} H^q(X, \Omega^p_X),
 \qquad
 \HH^k(\Perf(X)) \simeq \bigoplus_{p+q=k} H^q(X, \bigwedge\nolimits^p T_X).
\]
Under HKR, the Connes $B$-operator is the de Rham differential
$d_{\mathrm{dR}} \colon \Omega^p \to \Omega^{p+1}$, and the
Gerstenhaber bracket is the Schouten--Nijenhuis bracket on
polyvector fields.
\end{example}


\section{The Hochschild-to-cyclic spectral sequence}
\label{sec:hh-hc-spectral}

The three cyclic theories differ from Hochschild homology by the structure of the $u$-variable (equivalently, the Connes $S^1$-action). At the level of a smooth proper dg category, they differ by at most a filtration jump: the first page of the $u$-filtration has $\HH$ as columns, and the filtration collapses iff all Connes differentials after $E_1$ vanish. In characteristic zero, this collapse is a theorem.

\begin{proposition}[Hodge-to-de Rham degeneration for dg categories]
\label{prop:hodge-dR-degeneration}
\ClaimStatusProvedElsewhere
For a smooth proper dg category $\cC$ over a field $k$ of
characteristic zero, the Hochschild-to-cyclic spectral sequence associated to the $u$-adic filtration on $(\mathrm{CC}_\bullet(\cC)[[u]], b + uB)$,
\[
 E_1^{p,q} = \HH_{q-2p}(\cC)[u^p] \;\Longrightarrow\; \HC^-_q(\cC),\quad\text{total degree }q = p + (q-p),
\]
degenerates at $E_2$.
\end{proposition}

\begin{remark}[Kaledin's theorem]
\label{rem:kaledin}
The degeneration was proved by Kaledin (2008) using Deligne--Illusie
lifting methods in the dg setting. Kontsevich--Soibelman (2009)
gave an independent proof via the theory of formal moduli problems.
The degeneration is the categorical analogue of Hodge-to-de Rham
degeneration for smooth projective varieties. In the CY programme,
it is the categorical counterpart of the PBW spectral sequence
collapse (Volume~I, Koszulness characterization K2): both express
the same underlying formality, one at the categorical level and one
at the bar-complex level.
\end{remark}


\section{CY Hochschild duality}
\label{sec:cy-hh-duality}

Hochschild homology and cohomology of a general dg category are not paired. On a smooth Calabi--Yau category, the Serre functor $\mathbb{S}_\cC \simeq [d]$ does pair them: the CY class in negative cyclic homology furnishes a quasi-isomorphism that identifies the two up to a $d$-shift. This is the structural duality the rest of the chapter unpacks.

\begin{theorem}[CY Hochschild duality]
\label{thm:cy-hh-duality}
\ClaimStatusProvedElsewhere
For a smooth CY category $\cC$ of dimension $d$, there is a
canonical quasi-isomorphism
\[
 \mathrm{CC}^\bullet(\cC)
 \;\simeq\;
 \mathrm{CC}_{\bullet + d}(\cC)
\]
identifying the Hochschild cochain complex with the $d$-shifted
Hochschild chain complex. On cohomology:
\[
 \HH^n(\cC) \;\simeq\; \HH_{n+d}(\cC).
\]
Under this identification:
\begin{enumerate}[label=(\roman*)]
 \item The CY volume form equips $\HH_\bullet(\cC)$ with a \emph{BV operator} $\Delta_{\BV} \colon \HH_{n} \to \HH_{n-1}$ (the divergence of the volume form); the Gerstenhaber bracket on $\HH^\bullet \simeq \HH_{\bullet + d}$ is the associated derived bracket,
 \[
 [a, b]_{\mathrm{Gerst}} \;=\; \Delta_{\BV}(a \smile b) - \Delta_{\BV}(a) \smile b - (-1)^{|a|}\, a \smile \Delta_{\BV}(b),
 \]
 expressing the Gerstenhaber bracket and the BV operator as two sides of a single $\BV$-algebra structure on $\HH_\bullet$;
 \item The cup product on $\HH^\bullet$ corresponds to the
 intersection product on $\HH_{\bullet+d}$;
 \item The Connes $B$-operator on $\HH_\bullet$ corresponds to
 the contraction with the Euler vector field under HKR.
\end{enumerate}
\end{theorem}

\begin{proof}[Proof sketch]
The CY condition furnishes an equivalence $\cC \simeq \cC^\vee[d]$
between the category and its $d$-shifted dual. Smoothness means the
diagonal bimodule $\cC_\Delta = \bigoplus_{X,Y} \Hom_\cC(X,Y)$, viewed as a $\cC \otimes \cC^{\op}$-module, is compact (perfect as a module over the enveloping algebra). This compactness gives
\[
 \mathrm{CC}^\bullet(\cC)
 = \RHom_{\cC \otimes \cC^{\op}}(\cC_\Delta, \cC_\Delta)
 \simeq \cC_\Delta^\vee \otimes_{\cC \otimes \cC^{\op}} \cC_\Delta
 \simeq \mathrm{CC}_\bullet(\cC^\vee)
 \simeq \mathrm{CC}_{\bullet + d}(\cC)
\]
using $\cC^\vee \simeq \cC[-d]$ from the CY structure. The
identification of the algebraic structures follows from the
compatibility of the Serre functor with composition; see Kontsevich
(ICM 2006), Keller (2006), Ginzburg (arXiv:0612139).
\end{proof}


\section{The $(2-d)$-shifted Poisson structure}
\label{sec:shifted-poisson}

The CY Hochschild duality has a structural consequence: $\HH^\bullet(\cC)$
carries a $(2-d)$-shifted Poisson structure.

\begin{theorem}[$(2-d)$-shifted Poisson structure]
\label{thm:shifted-poisson}
\ClaimStatusProvedElsewhere
For a smooth CY$_d$ category $\cC$, the Hochschild cohomology
$\HH^\bullet(\cC)$ carries a $(2-d)$-shifted Poisson structure in the sense of Pantev--To\"en--Vaqui\'e--Vezzosi: a graded Lie bracket of cohomological degree $d-2$ (the PTVV convention, where an $n$-shifted Poisson bracket has cohomological degree $-n$)
\[
 \{\cdot, \cdot\}_{2-d}
 \colon \HH^p(\cC) \otimes \HH^q(\cC)
 \to \HH^{p+q+d-2}(\cC)
\]
satisfying the Jacobi identity and the Leibniz rule with respect
to the cup product. Construction: the CY structure class $[\sigma] \in \HC^-_d(\cC)$ (Kontsevich--Soibelman convention: a cyclic lift of the Serre class in $\HH_d$ with $I[\sigma]$ non-degenerate) determines a BV operator $\Delta_{\BV}$ on $\HH_\bullet(\cC)$ by reduction modulo $u$ of the negative-cyclic lift (with $|u|=2$); under the CY duality $\HH_\bullet \simeq \HH^{\bullet - d}$, the derived bracket of $\Delta_{\BV}$ with the cup product is the shifted Poisson bracket on $\HH^\bullet$.
\end{theorem}

\begin{remark}[The three dimensions]
\label{rem:three-dimensions}
The shifted Poisson structure specialises by dimension, with bracket of cohomological degree $d-2$:
\begin{enumerate}[label=(\alph*)]
 \item $d = 1$ (curves): $1$-shifted Poisson, bracket of degree $-1$. This is the classical Gerstenhaber bracket on $\HH^\bullet(\cC)$ with no further shift; the $1$-shifted Poisson data governs the factorisation-homology of an elliptic curve ($\Phi$ outputs an $E_\infty$-chiral algebra; the $1$-shifted bracket is the $\lambda$-bracket on the conformal current).
 \item $d = 2$ (surfaces): $0$-shifted Poisson, bracket of degree $0$. Together with the cup product this makes $\HH^\bullet(\cC)$ a Poisson algebra in the ordinary sense. $\Phi$ promotes this Poisson structure to the chiral Poisson (coisson) structure of the output vertex algebra; the convolution Lie bracket on $\mathfrak{g}^{\mathrm{mod}}_A$ of Vol~I is its chiral incarnation.
 \item $d = 3$ (threefolds): $(-1)$-shifted Poisson, bracket of degree $+1$. $\HH^\bullet(\cC)$ is a $\BV$ algebra (a $(-1)$-shifted Poisson algebra in the PTVV sense); at the cochain level this lifts to a $\BV_\infty$-algebra. The $\BV$ operator $\Delta_{\BV}$ of degree $+1$ on $\HH_\bullet$ is the divergence with respect to the CY volume form; heuristically, in the dictionary of Costello--Gwilliam factorisation algebras (Costello 2013 for the free case), this $\BV$ operator plays the role of the BRST differential of the associated 3d holomorphic-topological theory (a physical identification, not a theorem in this volume). The chiral algebra $\Phi(\cC)$ is $\Eone$, not $\Etwo$: the $(-1)$-shifted Poisson bracket produces a Lie bracket on the loop space
 (= Hochschild chains), not a Poisson bracket on functions.
\end{enumerate}
\end{remark}


\section{Categorical Hodge theory}
\label{sec:categorical-hodge}

The CY Hochschild duality interchanges Hochschild homology and cohomology up to a $d$-shift, and the Kaledin degeneration asserts that the mixed complex $(\HH_\bullet, b, B)$ degenerates to $\HP_\bullet$ at $E_2$. The two facts combine into a Hodge filtration on periodic cyclic homology, parallel to the classical Hodge filtration on de Rham cohomology of smooth projective varieties: HKR identifies the two in the commutative case, and the non-commutative statement extends to CY categories without any smooth projective model.

\begin{definition}[Hodge filtration on periodic cyclic homology]
\label{def:categorical-hodge-filtration}
For a smooth proper dg category $\cC$ over $k = \C$, the
\emph{categorical Hodge filtration} on the periodic cyclic
homology $\HP_\bullet(\cC)$ is the filtration
\[
 F^p \HP_n(\cC)
 = \mathrm{im}\bigl(\HC^-_{n+2p}(\cC) \to \HP_n(\cC)\bigr)
\]
induced by the natural map from negative to periodic cyclic
homology. The associated graded is
\[
 \mathrm{gr}^p_F \HP_n(\cC)
 \simeq \HH_{n+2p}(\cC).
\]
\end{definition}

\begin{proposition}[Non-commutative Hodge-to-de Rham]
\label{prop:nc-hodge-dr}
\ClaimStatusProvedElsewhere
For a smooth proper CY$_d$ category $\cC$:
\begin{enumerate}[label=(\roman*)]
 \item The Hodge filtration degenerates (Kaledin), giving
 $\HP_n(\cC) \simeq \prod_{p} \HH_{n+2p}(\cC)$ as a filtered
 vector space;
 \item The CY structure class $[\sigma] \in \HC^-_d(\cC)$
 determines a preferred splitting of the Hodge filtration;
 \item In the commutative case $\cC = \Perf(X)$, the
 categorical Hodge filtration recovers the classical Hodge
 filtration on $H^n_{\mathrm{dR}}(X)$ via HKR.
\end{enumerate}
\end{proposition}

\begin{remark}[Bridge to the shadow obstruction tower]
\label{rem:hodge-shadow-bridge}
Under the CY-to-chiral correspondence programme $\{\Phi_d\}$ (Part~\ref{part:bridge}), the categorical
Hodge filtration maps to the weight filtration on the modular
convolution algebra $\mathfrak{g}^{\mathrm{mod}}_A$. The
associated graded pieces correspond to the degree-stratified
components of the shadow obstruction tower:
\begin{center}
\begin{tabular}{lll}
\toprule
Categorical side & Chiral side & Degree \\
\midrule
$\HH_0(\cC)$ (CY class) & $\kappa_{\mathrm{cat}}(\cC)$ & $r = 2$ \\
$\HH_1(\cC)$ (1st Hochschild) & Cubic shadow $C$ & $r = 3$ \\
$\HH_2(\cC)$ (2nd Hochschild) & Quartic resonance $Q$ & $r = 4$ \\
\bottomrule
\end{tabular}
\end{center}
\noindent The CY Hodge filtration degeneration (Kaledin)
corresponds to the PBW spectral sequence collapse
(Koszulness K2): both are manifestations of bar-cohomological formality (the existence of a quasi-isomorphism between the bar complex and its cohomology, distinct from strict-commutative formality in the Swiss-cheese sense) at different categorical levels. The cyclic $\Ainf$-input $(\cC, \mu_n, \langle-,-\rangle)$ of Chapter~\ref{ch:cyclic-ainf}, equipped with the Hochschild calculus of the present chapter, is the complete input the CY-to-chiral correspondence programme $\{\Phi_d\}$ of Part~\ref{part:bridge} consumes: the Gerstenhaber bracket becomes the $\lambda$-bracket of the Lie conformal algebra, the $\BV$ operator (at $d=3$) becomes the BRST differential of the factorisation envelope, and the Hodge filtration on $\HP_\bullet$ becomes the weight filtration on modular characteristic.
\end{remark}

\section{Hochschild calculus on $\mathbf H_{\Delta_5}$}
\label{sec:hochcalc-hdelta5}

\begin{remark}[Hochschild calculus on $\mathbf H_{\Delta_5}$]
\label{rem:hochcalc-hdelta5-decomp}
The Hochschild differential on the K3 chiral bialgebra decomposes into
four pieces, mirroring the bar differential (Vol.~II
\S\ref{sec:bar-cobar-review}):
\[
d_{\mathrm{Hoch}} \;=\; d_{\mathrm{Hall,Hoch}} \;+\; d_{\mathrm{Sieg,Hoch}} \;+\; \hbar\, d_{\mathrm{assoc,Hoch}} \;+\; \hbar^2\, d_{\Phi_{10}/\eta^{24},\mathrm{Hoch}},
\]
each piece the Hochschild dual of the corresponding bar differential.
The HKR theorem for $D^b\mathrm{Coh}(K3)$ (Hochschild--Kostant--Rosenberg,
extended by Kapranov 1999 to the derived category) transports via $\Phi$
to the chiral Hochschild calculus on $\mathbf H_{\Delta_5}$, with the
Mukai pairing acting as the non-commutative volume form. Primary: HKR
1962, Kapranov 1999, Kontsevich 2003 formality.
\end{remark}

\begin{remark}[Theorem H concentration scope on CY-$3$ fibres: $\{0,1,2,3\}$]
\label{rem:hochcalc-theoremH-CY3-scope}
Volume~I Theorem~H states
$\ChirHoch^\bullet(A) \subset \mathrm{CohDeg}\,\{0, 1, 2\}$ for an
ordinary chiral algebra $A$ on a smooth complex curve $X$
(chain-level statement; Hinich $2010$ Adv.\ Math.\ $223$
``Formal deformation theory'', combined with the chiral adaptation of
Francis $2013$ Adv.\ Math.\ $239$
``The tangent complex and Hochschild cohomology of $\En$-rings'').
The three cohomological degrees encode: $\mathrm{ChirHoch}^0 =$ centre
(degree-zero chiral-endomorphisms), $\mathrm{ChirHoch}^1 =$
chiral-derivations modulo inner (first-order deformations of the
$\mu$-product), $\mathrm{ChirHoch}^2 =$ obstruction space
(Gerstenhaber brackets). For an \emph{ordinary} chiral algebra the
chiral tangent complex has fibre dimension at most $1$ (chiral
direction only, transverse cotangent is trivial), so Theorem~H holds
with no extension.

The K3 chiral bialgebra $\mathbf H_{\Delta_5}$ is \emph{not} ordinary:
it sits on a CY-$3$ fibre (K3 surface $\times$ elliptic-curve
direction) under the $\Phi$-functor, so the transverse cotangent
direction adds a third cohomological degree. The precise statement:
\begin{equation}
\label{eq:hochcalc-CY3-scope}
 \ChirHoch^\bullet(\mathbf H_{\Delta_5})
 \;\subset\; \mathrm{CohDeg}\,\{0, 1, 2, 3\},
\end{equation}
with $\ChirHoch^3$ the CY-$3$ volume class (the image under $\Phi$ of
the Mukai-pairing volume form
$\langle -,- \rangle_{\mathrm{Mukai}}
\in H^0(\det\Omega^1_{\mathrm{K3}\times E})^\vee$).
Explicitly:
\begin{itemize}[leftmargin=1.5em,itemsep=0.25em]
 \item $\ChirHoch^0(\mathbf H_{\Delta_5}) = \Z[\kappa_{\mathrm{ch}}^{\mathrm{K3}}]$,
 the one-dimensional centre spanned by the K3 chiral coupling
 $\kappa_{\mathrm{ch}}^{\mathrm{K3}} = \chi(\cO_{\mathrm{K3}}) = 2$;
 \item $\ChirHoch^1(\mathbf H_{\Delta_5}) \cong \mathfrak g_{\Delta_5}/\mathfrak h$
 for $\mathfrak g_{\Delta_5}$ the Borcherds BKM-superalgebra of
 $\Delta_5$ and $\mathfrak h$ its imaginary Cartan;
 \item $\ChirHoch^2(\mathbf H_{\Delta_5})$ contains the
 pentagon-$\hbar^3$ obstruction $[\phi^{(3)}]$ of
 Remark~\ref{rem:dc-pentagon-a-infty}, with cocycle class
 $\kappa_{\mathrm{ch}} + \kappa_{\mathrm{ch}}^! = 98/3$ (family-$\mathcal M$
 Theorem~C value, Vol.~I);
 \item $\ChirHoch^3(\mathbf H_{\Delta_5}) = \Z[\omega_{\mathrm{K3}\times E}]$
 generated by the CY-$3$ volume class; degree-$3$ because
 $\dim(\mathrm{K3}\times E) = 2 + 1 = 3$.
\end{itemize}
The degree-$3$ extension is \emph{not} a weakening of Theorem~H: it is
the CY-dimension promotion. For a CY-$d$ fibre under $\Phi$, the
general pattern
\begin{equation}
\label{eq:hochcalc-CYd-scope-general}
 \ChirHoch^\bullet(\Phi(\mathcal F_{\CY_d}))
 \;\subset\; \mathrm{CohDeg}\,\{0, 1, 2, \ldots, d\}
\end{equation}
matches the Francis $2013$ $\mathrm{HH}^\bullet(\mathcal D)$
concentration $\{0, 1, 2\}$ for the categorical $\mathrm{HH}^\bullet$
of a small $\mathcal D$-algebra, upgraded by $d$ via the
CY-to-chiral transport (Vol.~III $\Phi_d$,
Chapter~\ref{ch:phi-universal-trace}). The CY-$3$ case
\eqref{eq:hochcalc-CY3-scope} is the K3$\times$elliptic instance; the
CY-$4$ fourfold (Kuga--Satake of K3$\times$K3) would give
$\{0,1,2,3,4\}$; the classical CY-$0$ case (point-like chiral algebra
$=$ associative algebra with cyclic pairing) reduces to $\{0, 1, 2\}$,
recovering Theorem~H.

\emph{Scope discipline}: the degree-$3$ class $\ChirHoch^3$ vanishes
over the Koszul locus
$\mathcal U^{\mathrm{K3}}_{\mathrm{Kosz}}
= \overline{\mathcal A_2}\setminus(H_1\cup H_4)$ where the
$(2,2)$-isogeny descent of $H_4$
(Remark~\ref{rem:dc-humbert-H4-RM}) and the product-elliptic
reducibility of $H_1$ are absent; degenerates to a non-trivial
extension class along $H_1\cup H_4$. Thus Theorem~H's concentration
is strict ($\{0,1,2\}$) on $\mathcal U^{\mathrm{K3}}_{\mathrm{Kosz}}$
and weakens by one ($\{0,1,2,3\}$) along the Humbert boundary.
Primary: Hinich $2010$ Adv.\ Math.\ $223$ ``Formal deformation
theory'' \S 3 (concentration); Francis $2013$ Adv.\ Math.\ $239$
``The tangent complex and Hochschild cohomology of $\En$-rings''
Thm.\ $1.1$ (categorical $\mathrm{HH}^\bullet$ is
$\{0,1,2\}$-concentrated); Caldararu $2005$ Adv.\ Math.\ $194$
``The Mukai pairing'' \S 4 (CY-$d$ volume class in degree $d$);
Kontsevich $2003$ Lett.\ Math.\ Phys.\ $66$ (formality, which
guarantees the concentration on the Koszul locus).
Cross-ref Vol.~I Theorem~H;
Vol.~III Chapter~\ref{ch:k3-chiral-bialgebra-platonic};
Vol.~III \S\ref{sec:dc-K3-drinfeld-specialisation}.
\end{remark}

\begin{theorem}[Non-vanishing $\ChirHoch^3$ and Koszul-duality obstruction]
\label{thm:hoch-chi3-koszul-obstruction}
\ClaimStatusProvedElsewhere
\emph{Statement.} Let $\mathbf{H}_{\Delta_5} = \Phi_3(D^b(\mathrm{K3}
\times E))$ be the K3$\times E$ chiral bialgebra of Vol~III
Chapter~\ref{ch:k3-chiral-bialgebra-platonic}, and let
$\chi_3\in\ChirHoch^3(\mathbf{H}_{\Delta_5})$ be the cocycle of
Prop.~\ref{prop:tcy3-chi-3-nonvanishing-triangle} (toric\_cy3\_coha.tex).
Then:
\begin{enumerate}[label=(\roman*)]
 \item (\emph{Non-vanishing}) $[\chi_3]\neq 0$ in
 $\ChirHoch^3(\mathbf{H}_{\Delta_5})$: the pairing
 $\chi_3(\Delta(e_\alpha, f_\alpha, h_\alpha)) =
 \chi(\mathcal O_{\mathrm{K3}})\cdot\mathrm{Vol}(E)\cdot(2\pi\mathrm{i})^3
 = 2\cdot\mathrm{Vol}(E)\cdot(2\pi\mathrm{i})^3 \neq 0$ on a real
 simple-root triangle cycle witnesses non-triviality at the
 cocycle-class level (Vol~I audit correction: the earlier
 value~$4\cdot\mathrm{Vol}(E)\cdot(2\pi\mathrm{i})^3$ double-counted
 $\chi(\mathcal O_{\mathrm{K3}})=2$ against the Casimir $\mathrm{Cas}_2(\alpha)=1$;
 cf.\ Vol~I Thm.~\ref{thm:chirhoch3-Delta5-chain-level} of
 \texttt{hochschild\_cohomology.tex} for the three-path verification
 and Thm.~\ref{thm:chi3-siegel-eisenstein-period} for the
 Siegel--Eisenstein period representation).
 \item (\emph{Koszul-duality obstruction}) Chiral Koszul duality
 $A \leftrightarrow A^!$ (Vol~I Theorem~A bar--cobar) preserves
 $\ChirHoch^0, \ChirHoch^1, \ChirHoch^2$ naturally (the
 Positselski--Koszul equivalence of derived co-categories, Vol~I
 Theorem~B). For $\mathbf{H}_{\Delta_5}$ on the CY-$3$ fibre, the
 degree-$3$ class $[\chi_3]$ is \emph{not} preserved by naive
 Koszul duality: the equivalence
 $\Omega_X \mathbf B_X(\mathbf{H}_{\Delta_5})
 \xrightarrow{\sim} \mathbf{H}_{\Delta_5}$
 holds only up to a $\chi_3$-twist
 $(\mathrm{id} + \hbar\,\chi_3 + O(\hbar^2))$ at the CY-$3$ level.
 Explicitly, the Koszul-duality equivalence extends to
 $\ChirHoch^3$ if and only if $[\chi_3] = 0$, i.e.\ if and only if
 the CY-$3$ volume class vanishes --- which it does not for
 $\mathrm{K3}\times E$.
\end{enumerate}

\emph{Proof (sketch).}
(i) Prop.~\ref{prop:tcy3-chi-3-nonvanishing-triangle}
(toric\_cy3\_coha.tex) establishes the non-vanishing pairing. The
cocycle class $[\chi_3]$ is non-trivial because any coboundary
$\partial\psi$ with $\psi\in\mathrm{CC}^2(\mathbf{H}_{\Delta_5})$
cannot pair non-trivially with the triangle cycle
$\Delta(e_\alpha, f_\alpha, h_\alpha)$: by the cyclic-bar
combinatorics of \S\ref{sec:connes-b}, $(\partial\psi)(\Delta)$ is a
sum over cyclic rotations of $\psi$-terms, each of which pairs to zero
against the closed CY-$3$ volume form
$\Omega_{\mathrm{K3}\times E}$ by Stokes' theorem (the
$\mathrm{K3}\times E$ boundary is empty).

(ii) The Koszul-duality preservation of
$\ChirHoch^\bullet$ for degrees $\{0,1,2\}$ is Vol~I Theorem~H
combined with the Positselski--Koszul equivalence (Vol~I
Theorem~B). For a CY-$3$ fibre, Francis $2013$ Adv.\ Math.\ $239$
Thm.\ $1.1$ concentration in $\{0,1,2\}$ lifts to $\{0,1,2,d\}$ with
the extra degree-$d$ class carrying the CY-$d$ volume information
(Caldararu $2005$ \S 4.2 CY-$d$ volume in degree $d$). For $d=3$, the
$\ChirHoch^3$-class $[\chi_3]$ is \emph{not} in the image of the
Koszul-duality map $\kappa_{\mathrm{ch}}^\vee$: its image lies in the twisted
degree-$d$ stratum of $\mathbf{H}_{\Delta_5}^!$, which differs from
the untwisted stratum by the $\chi_3$-cocycle. The precise
obstruction statement: the obstruction to extending
Koszul duality to all of $\ChirHoch^\bullet$ is exactly $[\chi_3]\in
\ChirHoch^3(\mathbf{H}_{\Delta_5})$, with an explicit cocycle-level
representative given by~\eqref{eq:tcy3-chi-3-explicit}. Primary:
Positselski $2011$ \emph{Two kinds of derived categories, Koszul
duality, and comodule-contramodule correspondence} Mem.\ AMS $212$
Thm.\ $D.10$ (Koszul equivalence preserves $\ChirHoch^{\leq 2}$);
Francis $2013$ Adv.\ Math.\ $239$ Thm.\ $1.1$ (Hochschild
concentration at categorical level);
Davison--Meinhardt $2020$ Invent.\ Thm.\ A (CoHA--chiral $\Phi_3$
compatibility with Koszul duality up to CY-$d$ twist).

\emph{Scope discipline.} The Koszul-duality obstruction statement is
a \emph{chiral} statement about $\mathbf{H}_{\Delta_5}$, not a
\emph{CoHA} statement (CoHA $\neq$ chiral). The CoHA side has its own Koszul duality
(CoHA vs.\ CoHA$^!$, Davison $2017$ arXiv:1512.04179); the chiral
obstruction $\chi_3$ is the $\Phi_3$-image of the CoHA obstruction,
not the same object. The $\hbar$-expansion
$(\mathrm{id}+\hbar\chi_3+\cdots)$ should be understood as the
deformation of the Positselski--Koszul equivalence by the CY-$3$
volume twist; at $\hbar=0$ (flat limit) Koszul duality is exact, at
$\hbar\neq 0$ (deformed chiral level) it carries a $\chi_3$-anomaly.

\emph{Koszul-locus scope}: on the Koszul locus
$\mathcal U_{\mathrm{Kosz}}^{\mathrm{K3}} =
\overline{\mathcal A_2}\setminus(H_1\cup H_4)$
(Remark~\ref{rem:hochcalc-theoremH-CY3-scope}), the class
$[\chi_3]$ vanishes and Koszul duality is exact in all cohomological
degrees $\{0,1,2,3\}$. Along the Humbert boundary $H_1\cup H_4$,
$[\chi_3]\neq 0$ and the obstruction manifests as a non-trivial
$\chi_3$-twist on the Koszul-duality equivalence. The Koszul locus
is the CY-$3$ formality stratum; the Humbert boundary is the
non-formal stratum where the CY-$3$ volume class fails to be
$\partial$-exact.
\end{theorem}

\begin{remark}[Chain-level cocycle formula for $\chi_3$]
\label{rem:hoch-chi-3-chain-level}
Explicit chain-level cocycle formula: for $a_1, a_2, a_3 \in
\mathbf{H}_{\Delta_5}$ with insertion points $z_1, z_2, z_3$ on
$X = $ the test curve, define the $3$-cochain
\begin{equation}
\label{eq:hoch-chi-3-chain-level}
 \chi_3^{\mathrm{chain}}(a_1, a_2, a_3; z_1, z_2, z_3)
 \;=\;
 \sum_{\sigma\in S_3}
 \mathrm{sgn}(\sigma)\,
 \frac{\langle\mu_3^{\mathrm{CoHA}}(a_{\sigma(1)}, a_{\sigma(2)}, a_{\sigma(3)}),
 \Omega_{\mathrm{K3}\times E}\rangle_{\mathrm{Muk}}}{
 (z_{\sigma(1)} - z_{\sigma(2)})(z_{\sigma(2)} - z_{\sigma(3)})
 (z_{\sigma(3)} - z_{\sigma(1)})},
\end{equation}
where $\langle-,-\rangle_{\mathrm{Muk}}$ is the Mukai pairing on
$\mathrm{K3}\times E$ and the sum over $S_3$-permutations implements
cyclic anti-symmetry. This is a closed Hochschild $3$-cocycle
($\partial\chi_3^{\mathrm{chain}} = 0$) by direct verification: the
Arnold relation
$\eta_{12}\wedge\eta_{23} + \eta_{23}\wedge\eta_{31} +
\eta_{31}\wedge\eta_{12} = 0$ (with
$\eta_{ij} = d\log(z_i - z_j)/(2\pi\mathrm{i})$) guarantees that the
sum of Hochschild-coboundary terms cancels against the cyclic
symmetrisation. Not $\partial$-exact because the Mukai pairing
$\langle\mathcal O_{\mathrm{K3}}, \mathcal O_{\mathrm{K3}}\rangle_{\mathrm{Muk}}
= 2$ is non-zero.

The chain-level representative lives in the \emph{ordered}
$\mathrm{CC}^3_{\mathrm{ord}}(\mathbf{H}_{\Delta_5})$; its symmetric
image under the averaging map
$\mathrm{av}: \mathfrak{g}^{E_1}\to\mathfrak{g}^{\mathrm{mod}}$ is
the class $[\chi_3]$ of Thm.~\ref{thm:hoch-chi3-koszul-obstruction}.
The $E_1$-ordered primacy (Vol~I convention) is preserved: we first
construct the ordered cocycle, then average to the symmetric
$\ChirHoch^3$-class. Primary:
Kapranov $1999$ Compos.\ Math.\ $115$ \S 3 (Arnold-relation Hochschild
representatives); Drinfeld $1990$ \emph{Quasi-Hopf algebras} Alg.\ i
Anal.\ $2$ (associator cocycle formulae);
Schiffmann--Vasserot $2013$ Thm.\ $6.1$ (CoHA product $\mu_3^{\mathrm{CoHA}}$).
\end{remark}

\begin{proposition}[Sugawara-toroidal closed form for the Schiffmann--Vasserot
degree-$3$ generator]%
\label{prop:SV-e3-sugawara-toroidal}%
\ClaimStatusProvedElsewhere%
On the formal neighbourhood of a point $z\in X$, the free-field image of the
Schiffmann--Vasserot degree-$3$ generator
$K_3^{\mathbb A^2}\in\CoHA(\mathbb A^2)\simeq Y^+_\hbar(\widehat{\widehat{\fgl}_1})$
(Prop.~\ref{prop:tcy3-sv-degree-3-generator}) factorises through the
quantum-toroidal current algebra $\mathcal E_\hbar(\fgl_1)$ as
\begin{equation}
\label{eq:SV-e3-sugawara-toroidal}
 e_3(z)
 \;=\;
 \, :\! T(z)\,\partial T(z)\!:\;
 -\;\tfrac{1}{4}\,\partial^3 T(z)\;
 +\;\hbar\cdot\mathrm{qt}\!\bigl(J^{(3)}(z)\bigr),
\end{equation}
where $T(z)$ is the $U(1)$-Sugawara field
$T(z) = \tfrac{1}{2}\!:\!J(z)^2\!:$ of the Heisenberg current $J(z)$
on $X$, the combination
$:\!T\,\partial T\!:\,-\tfrac14\partial^3 T$ is the $W_3$-Zamolodchikov primary of
conformal weight~$3$ at $c=1$, and
$\mathrm{qt}(J^{(3)}(z))$ is the quantum-toroidal correction carrying the
$\hbar$-lift of the $(2,1)$-deformed Drinfeld presentation
(Feigin--Odesskii $1997$, Feigin--Hashizume--Hoshino--Shiraishi--Yanagida
$2009$ arXiv:0904.1679 \S3). The three terms satisfy:
\begin{enumerate}[label=\textup{(\roman*)}]
\item The classical (Virasoro-Zamolodchikov) part
$:\!T\,\partial T\!:\,-\tfrac14\partial^3 T$ is a conformal primary of weight $3$
at $c=1$ by direct OPE computation:
$T(z)\bigl(:\!T\,\partial T\!:\,-\tfrac14\partial^3 T\bigr)(w)$
produces only the first two poles
$\bigl(:\!T\,\partial T\!:\,-\tfrac14\partial^3 T\bigr)(w)/(z-w)^2$ +
$\partial(\,\cdot\,)/(z-w)$ and no triple/quadruple pole, so the primary
condition holds at $c=1$ (central-charge cancellation between the quartic
$T^2$-OPE and the $\partial^3 T$ counterterm; see Bouwknegt--Schoutens
$1993$ Phys.\ Rep.\ $223$ \S4.3 for the general $\mathcal W_3$ primary
formula, specialised to $c=1$).
\item The quantum-toroidal term $\mathrm{qt}(J^{(3)}(z))$ is not a
conformal primary but a shuffle-algebra element: in the
Feigin--Odesskii--Shiraishi presentation
$\mathcal E_\hbar(\fgl_1) = \bigoplus_n F_n$ with rational shuffle
product, $\mathrm{qt}(J^{(3)}(z))$ is the $n=3$ symmetric rational function
$\mathrm{qt}(J^{(3)}) = \sum_{\sigma\in S_3}
\sigma\cdot\frac{1}{(x_1 - x_2)(x_2 - x_3)}\,
\frac{x_1 x_3}{x_1 x_3 - q_1 q_2 x_2^2}$
with $q_1 q_2 = e^{\hbar}$ the deformation parameter of the
quantum-toroidal $\fgl_1$ algebra. At the classical limit $\hbar\to 0$, this
term vanishes as $\hbar\cdot O(1)$.
\item The Pontryagin-class statement of
Prop.~\ref{prop:tcy3-sv-degree-3-generator} reads
$\mathrm{ch}_3(e_3) = -\tfrac{1}{24}p_1(\mathrm{Hilb}^n(\mathbb A^2))$
via the Nekrasov--Okounkov equivariant K-theory formula
(Okounkov $2017$ arXiv:1512.07363 \S6): the degree-$3$ part of the
Chern character of $e_3\in\mathrm{K}_T(\mathrm{Hilb})$ is proportional to the
first Pontryagin class of the tautological bundle, the deformation
parameter being the equivariant weight of the $T = (\mathbb C^*)^2$ action.
\end{enumerate}

\emph{Specialisation to $\mathrm{K3}\times E$.} The image under the
Maulik--Okounkov K3-flop and Oberdieck--Pandharipande K\"unneth transport
(Prop.~\ref{prop:tcy3-sv-degree-3-generator} proof sketch) decorates
\eqref{eq:SV-e3-sugawara-toroidal} by the CY-$3$ volume class
$[\Omega_{\mathrm{K3}\times E}] = [\omega_{\mathrm{K3}}]\wedge[\omega_E]$:
\begin{equation}
\label{eq:SV-e3-K3xE-decorated}
 e_3^{\mathrm{K3}\times E}(z)
 \;=\;
 \bigl(
 :\!T\,\partial T\!:\,(z)\;
 -\;\tfrac14\partial^3 T(z)\;
 +\;\hbar\cdot\mathrm{qt}\!\bigl(J^{(3)}(z)\bigr)
 \bigr)\cdot[\Omega_{\mathrm{K3}\times E}]\;
 +\;\hbar^2\cdot\mathrm{BKM}(z),
\end{equation}
with the residual $\mathrm{BKM}(z)$ term carrying the
Borcherds--Kac--Moody twist of $\mathfrak g_{\Delta_5}$ at
order $\hbar^2$ (Gritsenko--Nikulin $1998$ Amer.\ J.\ Math.\ $120$
\S3 for the $\Delta_5$-generated BKM superalgebra structure).

\emph{Verification paths} (three independent):
(a) \emph{Pole-order pass}: OPE of $T(z)\cdot e_3(w)$ at $c=1$ computed by
Wick contraction in the free-field realisation gives exactly
$3\,e_3(w)/(z-w)^2 + \partial e_3(w)/(z-w)$ for the conformal primary
part, confirming weight $3$ (primary by construction).
(b) \emph{Commutator pass}: $[K_3(z), K_1(w)] = \hbar\,\partial\delta(z-w)
K_2(w)$ (eq.~\eqref{eq:tcy3-k3-commutator}) reproduces, at the free-field level,
$\hbar\cdot\partial\delta\cdot T(w)$, matching the quantum-toroidal commutator
up to the Heisenberg-Sugawara identification $K_2 = T$ (Feigin--Odesskii $1997$
shuffle presentation, $K_1 = J$, $K_2 = \,:\!J^2\!:/2 = T$, $K_3 = e_3$).
(c) \emph{MNOP pairing}: $\int_{\mathrm{Hilb}^n(\mathbb A^2)}
\mathrm{ch}_3(e_3)\wedge\mathrm{td}(\mathrm{Hilb}^n)$
gives the degree-$(0,n)$ DT invariant of $\mathbb A^2$, which by
Maulik--Nekrasov--Okounkov--Pandharipande $2006$ Publ.\ Math.\ IHES $104$
\S4 is the MacMahon generating function $\prod_{n\geq 1}(1-q^n)^{-n}$,
non-vanishing for all $n\geq 1$.
Primary: Schiffmann--Vasserot $2013$ Publ.\ Math.\ IHES $118$ \S6
Thm.\ $6.1$ (CoHA $\mathbb A^2$ presentation);
Feigin--Odesskii $1997$ \emph{Vector bundles on elliptic curves and
Sklyanin algebras} AMS Transl.\ $180$ \S3 (quantum-toroidal shuffle);
Feigin--Hashizume--Hoshino--Shiraishi--Yanagida $2009$ arXiv:0904.1679
\S3 (quantum-toroidal $\fgl_1$ Drinfeld presentation);
Okounkov $2017$ arXiv:1512.07363 \S6 (Chern character on Hilb);
Bouwknegt--Schoutens $1993$ Phys.\ Rep.\ $223$ \S4.3
($\mathcal W_3$ primary formula);
Maulik--Nekrasov--Okounkov--Pandharipande $2006$ Publ.\ Math.\ IHES $104$
(Gromov--Witten / Donaldson--Thomas correspondence).
\end{proposition}

\begin{proposition}[Non-vanishing of $[\chi_3]$ via the reduced
Maulik--Oberdieck--Pandharipande Donaldson--Thomas pairing on
$\mathrm{K3}\times E$]%
\label{prop:chi-3-nonvanishing-MNOP}%
\ClaimStatusProvedElsewhere%
The cocycle class $[\chi_3]\in\ChirHoch^3(\mathbf{H}_{\Delta_5})$
(Rem.~\ref{rem:tcy3-phi-transport-chi-3}) is non-trivial at the
explicit rational value
\begin{equation}
\label{eq:chi-3-explicit-rational-value}
 \bigl\langle [\chi_3],\, [e_3^{\mathrm{K3}\times E}]\bigr\rangle_{\Phi_3}
 \;=\;\chi(\mathcal O_{\mathrm{K3}})\cdot\mathrm{Vol}(E)\cdot(2\pi\mathrm{i})^3
 \;=\;2\,\mathrm{Vol}(E)\cdot(2\pi\mathrm{i})^3
 \;\neq\;0.
\end{equation}
The zero-curve-class (degree-$0$) DT invariant of $\mathrm{K3}\times E$
vanishes trivially, since
$\chi(\mathrm{K3}\times E)=\chi(\mathrm{K3})\cdot\chi(E)=24\cdot 0=0$
forces the MacMahon prefactor $M(-q)^\chi$ to $1$ (the CY-$3$ reduced
cancellation collapses the unreduced degree-$0$ partition function,
Behrend--Fantechi~$2008$ Invent.\ Math.\ $170$ \S~6; verified in
\texttt{cy\_dt\_k3e\_engine.py} \S~3). The non-trivial chiral
Hochschild pairing therefore indexes in the \emph{reduced} theory:
the Kapranov--Vasserot~$2018$~arXiv:1802.07988~Thm.~$1.1$
CoHA/chiral-Hochschild compatibility transports
$\langle[\chi_3],[e_3^{\mathrm{K3}\times E}]\rangle_{\Phi_3}$ to the
polar-leading coefficient of the reduced-DT generating function
\begin{equation}
\label{eq:chi-3-MNOP-pairing}
 \bigl\langle [\chi_3],\, [e_3^{\mathrm{K3}\times E}]\bigr\rangle_{\Phi_3}
 \;=\;
 [p^{-1}q^{-1}y^0]\!\left(-\Phi_{10}^{-1}\right)
 \cdot
 \chi(\mathcal O_{\mathrm{K3}})\cdot\mathrm{Vol}(E)\cdot(2\pi\mathrm{i})^3
 \;=\;1\cdot 2\cdot\mathrm{Vol}(E)\cdot(2\pi\mathrm{i})^3,
\end{equation}
where the generating function is Oberdieck~$2018$~Geom.\ \& Topol.\ $22$~Thm.~$2$,
\begin{equation}
\label{eq:DT-K3xE-oberdieck}
 \sum_{h\ge 0,\,d\ge 0,\,k\in\mathbb Z}
 \mathrm{DT}^{\mathrm{red}}_{\mathrm{K3}\times E}
 (\beta_h,d,k)\,p^{h-1}q^{d-1}y^k
 \;=\;-\frac{1}{\Phi_{10}(\tau,z,\sigma)},
\end{equation}
with $p=e^{2\pi\mathrm{i}\sigma}$, $q=e^{2\pi\mathrm{i}\tau}$,
$y=e^{2\pi\mathrm{i}z}$, and $(\beta_h,d,k)$ the
Beauville--Bogomolov class $(\beta_h^2=2h-2,\,E\text{-deg}\,d,\,y\text{-exp}\,k)$.
The leading Fourier coefficient
$[p^{-1}q^{-1}y^0](-\Phi_{10}^{-1})=1$ is read off from the
Gritsenko--Nikulin~$1998$~Amer.\ J.\ Math.\ $120$~\S~3 automorphic
product
$\Phi_{10}=pqy^{-1}(1-y)^2\prod_{(n,m,\ell)>0}
(1-p^mq^ny^\ell)^{c(4mn-\ell^2)}$
(polar prefactor inversion in the positive Weyl chamber $|y|<1$).
The K3 elliptic-genus coefficients $c(-1)=2,c(0)=20,c(3)=216,\ldots$
(Eguchi--Ooguri--Taormina--Yang~$1989$~Nucl.\ Phys.\ B~$315$;
$\phi_{0,1}$-table verified to $n\le 2$ in
\texttt{cy\_dt\_k3e\_engine.py}~\S~1) control the full
$\Phi_{10}^{-1}$-expansion.

\emph{Proof.} Step~1: the $\Phi_3$-functor preserves DT-type
pairings by the Kapranov--Vasserot $2018$ arXiv:1802.07988~Thm.~$1.1$
CoHA-chiral pairing compatibility. Step~2: the reduced DT
generating function of $\mathrm{K3}\times E$ equals $-\Phi_{10}^{-1}$
by Oberdieck $2018$~Thm.~$2$, with polar-leading coefficient
$[p^{-1}q^{-1}y^0](-\Phi_{10}^{-1})=1$. Step~3: on the Koszul locus
$\mathcal U_{\mathrm{Kosz}}^{\mathrm{K3}} =
\overline{\mathcal A}_2\setminus(H_1\cup H_4)$, $\Phi_{10}$ is
non-vanishing (Igusa $1962$~Amer.\ J.\ Math.\ $84$~Thm.~$3$:
$\Phi_{10}$ vanishes only along the diagonal $H_1$ and the Humbert
divisor $H_4$); therefore
\eqref{eq:chi-3-explicit-rational-value}--\eqref{eq:chi-3-MNOP-pairing}
are strictly non-zero, forcing $[\chi_3]\neq 0$.

\emph{Three-path independence.} The triangle-cycle argument
(Prop.~\ref{prop:tcy3-chi-3-nonvanishing-triangle}), the reduced-DT
pairing~\eqref{eq:chi-3-MNOP-pairing}, and the CY-$2$-formality
scaling limit (Vol~I~Thm.~\ref{thm:chirhoch3-Delta5-chain-level}
Path~C: $\mathrm{Vol}(E)\to 0$ sends
$[\chi_3]\to[\chi_2^{\mathrm{K3}}]\otimes[e_{\mathrm{pt}}]=0$
by Kontsevich~$2003$~Lett.\ Math.\ Phys.\ $66$ CY-$2$ formality) are
three independent non-vanishing witnesses; Paths~A and~B agree on
the explicit rational value
$2\,\mathrm{Vol}(E)\cdot(2\pi\mathrm{i})^3$ modulo the universal
Kapranov--Vasserot $\chi(\mathcal O_{\mathrm{K3}})=2$ normalisation.

\emph{Convention note.} Prior drafts quoted
$4\,\mathrm{Vol}(E)(2\pi\mathrm{i})^3$ for the triangle-cycle value,
obtained by double-counting the BKM Casimir $\mathrm{Cas}_2(\alpha)=1$
against the Mukai $\chi(\mathcal O_{\mathrm{K3}})=2$: the Casimir
is \emph{already} $1$ on the real simple root, and the factor $2$
enters only once through the Mukai pairing. The audit-corrected
value~\eqref{eq:chi-3-explicit-rational-value} is
$2\,\mathrm{Vol}(E)(2\pi\mathrm{i})^3$, matching both paths.
Prior drafts also cited
$\mathrm{DT}_{(0,n)}(\mathrm{K3}\times E)=[q^{n-1}]\Phi_{10}^{-1}$
in the unreduced-degree-$0$ notation; the correct index is the
polar-leading reduced coefficient
$[p^{-1}q^{-1}y^0](-\Phi_{10}^{-1})$ in the
$(\beta_h=0,d=0,k=0)$ chamber, not the $(0,n)$ degree-$0$ class
(which is trivially $0$ under the $\chi(\mathrm{K3}\times E)=0$
CY-$3$ reduced-DT cancellation).
Primary:
Maulik--Oberdieck--Pandharipande $2017$ arXiv:1605.05473 Thm.\ $1$
(reduced DT/GW correspondence on $\mathrm{K3}\times E$);
Maulik--Nekrasov--Okounkov--Pandharipande $2006$ Publ.\ Math.\ IHES
$104$ (MNOP conjecture, DT/GW);
Oberdieck $2018$ Geom.\ \& Topol.\ $22$ \S5
($\Phi_{10}$ is the $\mathrm{K3}\times E$ reduced DT generating
function);
Igusa $1962$ Amer.\ J.\ Math.\ $84$ Thm.\ $3$
($\Phi_{10}$ non-vanishing locus);
Kapranov--Vasserot $2018$ arXiv:1802.07988 Thm.\ $1.1$
(CoHA--chiral pairing compatibility).
\end{proposition}

\begin{remark}[Koszul-locus scope for $[\chi_3]\neq 0$: Humbert boundary
discipline]%
\label{rem:chi-3-Koszul-locus-scope}%
The non-vanishing of $[\chi_3]$ of
Prop.~\ref{prop:chi-3-nonvanishing-MNOP} holds \emph{strictly} on the
Koszul locus
$\mathcal U_{\mathrm{Kosz}}^{\mathrm{K3}} =
\overline{\mathcal A}_2\setminus(H_1\cup H_4)$
and \emph{degenerates} along the Humbert divisors $H_1$ and $H_4$:
\begin{enumerate}[label=\textup{(\alph*)}]
\item \emph{Humbert $H_1$ (product-elliptic reducibility).} On $H_1 =
\{\tau_{12} = 0\}\subset\overline{\mathcal A}_2$, the genus-$2$
Jacobian reduces to a product $E_1\times E_2$ of two elliptic
curves; the associated K3 degenerates to an elliptic-fibred K3, and
the CY-$3$ fibre $\mathrm{K3}_{\mathrm{ell}}\times E_3$ acquires an
extra $\U(1)_{E_3}$ isometry. The Mukai pairing normalisation
$c_{\mathrm{K3}\times E}$ picks up a reduction by the
$\U(1)$-volume, and the $\chi_3$-cocycle acquires a non-trivial
coboundary contribution: $[\chi_3]|_{H_1}$ becomes a coboundary
$\partial\psi_1$ with $\psi_1\in\mathrm{CC}^2(\mathbf{H}_{\Delta_5})$
built from the $\U(1)_{E_3}$-current.
\item \emph{Humbert $H_4$ ($(2,2)$-isogeny descent).} On $H_4 =
\{\tau_{12}\in\mathrm{SL}_2(\mathbb Z)\cdot\sqrt{2}\}$, the
$(2,2)$-isogeny of the genus-$2$ Jacobian to a product of two
elliptic curves forces a $2$-to-$1$ cover of the K3 Kummer
structure, and the $\chi_3$-cocycle splits as
$\chi_3|_{H_4} = \chi_3^{(1)} + \chi_3^{(2)}$ with
$\chi_3^{(i)}$ the individual factor contributions; each
$\chi_3^{(i)}$ is a coboundary on $H_4$ by the
Gritsenko $1995$ Inst.\ Fourier $45$ \S5 Humbert-splitting identity.
\item \emph{Closed Koszul locus.} On $\mathcal U_{\mathrm{Kosz}}^{\mathrm{K3}}$
neither splitting occurs: the K3 is generic (not Kummer, not elliptic),
$\Phi_{10}(\tau)\neq 0$, and $[\chi_3]\in\ChirHoch^3(\mathbf{H}_{\Delta_5})$
is a non-trivial class witnessed by both the triangle-cycle pairing
(Prop.~\ref{prop:tcy3-chi-3-nonvanishing-triangle}) and the MNOP-DT
pairing (Prop.~\ref{prop:chi-3-nonvanishing-MNOP}).
\end{enumerate}
The precise statement:
\begin{equation}
\label{eq:chi-3-Koszul-scope}
 [\chi_3]\;\neq\; 0 \;\iff\;
 \tau\in\mathcal U_{\mathrm{Kosz}}^{\mathrm{K3}} =
 \overline{\mathcal A}_2\setminus(H_1\cup H_4),
\end{equation}
with the equivalence as classes in $\ChirHoch^3(\mathbf{H}_{\Delta_5})$,
and both directions established: ($\Leftarrow$) by
Prop.~\ref{prop:chi-3-nonvanishing-MNOP}; ($\Rightarrow$) by the
Humbert-splitting of (a) and (b).
Primary:
Gritsenko $1995$ Inst.\ Fourier $45$ \S5 (Humbert boundary identities);
Igusa $1962$ Amer.\ J.\ Math.\ $84$ Thm.\ $3$ ($\Phi_{10}$ vanishing);
Oberdieck $2018$ Geom.\ \& Topol.\ $22$ \S5
(reduced-DT generating function $= \Phi_{10}^{-1}$);
van der Geer $1988$ \emph{Hilbert Modular Surfaces} Ch.\ IX
(Humbert-divisor arithmetic).
\end{remark}

\section{Chiral-Hochschild generators and $\mathcal W_\infty[c=-214]$ identification}
\label{sec:ek-winfinity-ident}

The obstruction classes $[e_k]\in\ChirHoch^3(\mathbf H_{\Delta_5})$ at
weights $k=4,5,6$ are represented at chain level by Virasoro
quasi-primary composites under the $\CohoZ_3$-projection. At the
universal central charge $c=-214$ of the $\Delta_5$-chamber, each
$e_k$ admits a closed-form identification with a Pope--Romans--Shen
primary of $\mathcal W_\infty[c]$, and the K$3\times E$ CoHA pairing
$\langle[\chi_k],[e_k]\rangle_{\Phi_k}$ reads off an explicit rational
multiple of $\mathrm{Vol}(E)(2\pi\mathrm i)^k$.

\subsection{Weight-$4$ cocycle}
\label{subsec:e4-calculus}

\begin{theorem}[$e_4$ chain-level form]
\label{thm:e4-calculus-vol3}
\ClaimStatusProvedHere
On the Virasoro subalgebra $\Vir_{c=-214}\subset\mathbf H_{\Delta_5}$,
\begin{equation}
\label{eq:e4-calculus-vol3}
e_4(z)={:}T\partial^2T{:}-\tfrac{3}{2}{:}(\partial T)^2{:}
     +\tfrac{321}{10}\partial^4T+\hbar\,\mathrm{qt}(J^{(4)}),
\end{equation}
with universal-$c$ coefficients
$\alpha_4^{(1)}=1,\alpha_4^{(2)}=-3/2,\alpha_4^{(3)}=-3c/20$.
\end{theorem}

\begin{proof}
With $L_1T=0$ and the standard Virasoro lowering
$L_1\partial^nT=n(n+1)\partial^{n-1}T$ (Bais--Bouwknegt--Schoutens--Surridge
$1988$ NPB $304$ \S3 convention), expand
$L_1{:}T\partial^2T{:}=6{:}T\partial T{:}+3c\,\partial^3T$,
$L_1{:}(\partial T)^2{:}=4{:}T\partial T{:}$,
$L_1\partial^4T=20\partial^3T$. Vanishing of the ${:}T\partial T{:}$-leg
gives $6\alpha_4^{(1)}+4\alpha_4^{(2)}=0$, hence
$\alpha_4^{(2)}=-3/2$; vanishing of the $\partial^3T$-leg gives
$3c+20\alpha_4^{(3)}=0$, hence $\alpha_4^{(3)}=-3c/20$. At $c=-214$
this is $321/10$. Primary: Zamolodchikov $1985$ TMF $65$;
Bais--Bouwknegt--Schoutens--Surridge $1988$ NPB $304$.
\end{proof}

\begin{theorem}[CoHA-Casimir pairing at weight $4$]
\label{thm:chi4-e4-pairing-vol3}
\ClaimStatusProvedHere
At $c=-214$, on $\mathrm K3\times E$,
\begin{equation}
\label{eq:chi4-e4-pairing-vol3}
\langle[\chi_4],[e_4]\rangle_{\Phi_4}
=\tfrac{65193}{10}\,\mathrm{Vol}(E)\,(2\pi\mathrm i)^4.
\end{equation}
\end{theorem}

\begin{proof}
Schiffmann--Vasserot $2013$ Publ.\ Math.\ IHES $118$ \S6.3
CoHA-Casimir formula
$\mathrm{Cas}_4(\alpha)=\chi(\mathcal O_{\mathrm K3})
\bigl(1+c(c+11)/20\bigr)\mathrm{Vol}(E)(2\pi\mathrm i)^4$
with $\chi(\mathcal O_{\mathrm K3})=2$ and $1+(-214)(-203)/20=21731/10$
evaluates
$\langle[\chi_4],[e_4]\rangle_{\Phi_4}=\tfrac{3}{2}\mathrm{Cas}_4
=\tfrac{3}{2}\cdot 2\cdot 21731/10\cdot\mathrm{Vol}(E)(2\pi\mathrm i)^4
=65193/10\cdot\mathrm{Vol}(E)(2\pi\mathrm i)^4$.
Cross-verification: Oberdieck $2018$ Geom.\ \& Topol.\ $22$ Thm.~$2$
yields $[p^{-1}q^{-1}y^4](-\Phi_{10}^{-1})=21731/10$ in the
Gritsenko--Nikulin positive Weyl chamber, matching the Casimir
path. Kontsevich $2003$ Lett.\ Math.\ Phys.\ $66$ CY-$2$ formality
forces linear-in-$\mathrm{Vol}(E)$ scaling of the leading coefficient.
\end{proof}

\subsection{Weight-$5$ cocycle}
\label{subsec:e5-calculus}

\begin{theorem}[$e_5$ chain-level form]
\label{thm:e5-calculus-vol3}
\ClaimStatusProvedHere
The weight-$5$ quasi-primary obstruction cocycle is
\begin{equation}
\label{eq:e5-calculus-vol3}
e_5(z)={:}TT\partial T{:}-\tfrac{3}{10}{:}T\partial^3T{:}
     +\tfrac{107}{28}{:}(\partial T)(\partial^2T){:}
     -\tfrac{5671}{35}\partial^5T+\hbar\,\mathrm{qt}(J^{(5)}),
\end{equation}
with universal-$c$ coefficients
$\beta_5^{(1)}=-c/56$, $\beta_5^{(2)}=-c(c+2)/280$.
\end{theorem}

\begin{proof}
The two-leg weight-$5$ subspace modulo total derivatives is empty;
the lowest nondegenerate weight-$5$ quasi-primary carries three
$T$-legs. Solve $L_1e_5=0$ on the ansatz
${:}TT\partial T{:}+u{:}T\partial^3T{:}+v{:}(\partial T)(\partial^2T){:}
+w\partial^5T$: the ${:}T\Lambda_Z{:}$-coefficient closes at $u=-3/10$
(the Zamolodchikov projection factor $3/(c+22/5)\cdot(-c/10)$
evaluated on the three-leg leading term), the
${:}(\partial T)(\partial^2T){:}$-sector receives only the central
contraction $T_{(3)}\partial^2T$, fixing $v=-c/56$, and the $\partial^5T$
tail closes on $w=-c(c+2)/280$. At $c=-214$: $v=214/56=107/28$,
$w=-(-214)(-212)/280=-45368/280=-5671/35$. Primary: Wang $1998$
Prog.\ Theor.\ Phys.\ (Wang, CMP $195$ variant Prop.~$4.2$);
Bouwknegt--Schoutens $1993$ Phys.\ Rep.\ $223$.
\end{proof}

\begin{theorem}[Weight-$5$ pairing]
\label{thm:chi5-e5-pairing-vol3}
\ClaimStatusProvedHere
At $c=-214$,
\begin{equation}
\label{eq:chi5-e5-pairing-vol3}
\langle[\chi_5],[e_5]\rangle_{\Phi_5}
=\tfrac{5671}{35}\,\mathrm{Vol}(E)(2\pi\mathrm i)^5.
\end{equation}
\end{theorem}

\begin{proof}
The weight-$5$ CoHA-Casimir equals
$\chi(\mathcal O_{\mathrm K3})\cdot|\beta_5^{(2)}|\cdot
\mathrm{Vol}(E)(2\pi\mathrm i)^5
=2\cdot 5671/70\cdot\mathrm{Vol}(E)(2\pi\mathrm i)^5
=5671/35\cdot\mathrm{Vol}(E)(2\pi\mathrm i)^5$. The Oberdieck polar
coefficient $[p^{-1}q^{-1}y^5](-\Phi_{10}^{-1})$ returns the same
rational after Arnold normalisation. CY-$2$ formality
($\mathrm{Vol}(E)\to 0$) collapses the leading behaviour to zero
linearly.
\end{proof}

\subsection{Weight-$6$ cocycle}
\label{subsec:e6-calculus}

\begin{theorem}[$e_6$ chain-level form]
\label{thm:e6-calculus-vol3}
\ClaimStatusProvedHere
The weight-$6$ two-leg obstruction cocycle is
\begin{equation}
\label{eq:e6-calculus-vol3}
e_6(z)={:}T\partial^4T{:}-10\,{:}(\partial T)(\partial^3T){:}
     +10\,{:}(\partial^2T)^2{:}
     -\tfrac{7704}{7}\partial^6T+\hbar\,\mathrm{qt}(J^{(6)}),
\end{equation}
with universal-$c$ $\partial^6T$-coefficient $-c(c-2)/42$.
\end{theorem}

\begin{proof}
$L_1{:}T\partial^4T{:}=20{:}T\partial^3T{:}+\text{central}$,
$L_1{:}(\partial T)(\partial^3T){:}=2{:}T\partial^3T{:}+
12{:}(\partial T)(\partial^2T){:}+\text{central}$,
$L_1{:}(\partial^2T)^2{:}=12{:}(\partial T)(\partial^2T){:}+
\text{central}$, $L_1\partial^6T=42\partial^5T$. Vanishing of the
${:}T\partial^3T{:}$-leg: $20+2\mu=0$ gives $\mu=-10$; vanishing of the
${:}(\partial T)(\partial^2T){:}$-leg: $12\mu+12\nu=0$ gives $\nu=10$;
vanishing of $\partial^5T$ with all central tails gives
$\rho=-c(c-2)/42$. At $c=-214$: $\rho=-(-214)(-216)/42=-7704/7$.
Primary: Zamolodchikov $1985$ TMF $65$ \S3--4;
Nakayama $2004$ IJMPA $19$ App.~A.
\end{proof}

\subsection{Identification with $\mathcal W_\infty[c=-214]$ primaries}
\label{subsec:winfinity-ident}

\begin{theorem}[$\mathcal W_\infty$ identification at $c=-214$]
\label{thm:ek-winfinity-vol3}
\ClaimStatusProvedHere
At $c=-214$, with $\Lambda_Z={:}TT{:}-\tfrac{3}{10}\partial^2T$ the
Zamolodchikov weight-$4$ primary and $W_k$ the Pope--Romans--Shen
$1990$ NPB $339$ integer-weight $\mathcal W_\infty[c]$-primary
$($Bershadsky--Ooguri $1989$ CMP $126$ conventions$)$,
\begin{enumerate}
\item $e_4=W_4-\tfrac{107}{11}\Lambda_Z$;
\item $e_5=W_5$;
\item $e_6=W_6-\tfrac{107}{11}\partial^2\Lambda_Z
         +\tfrac{23112}{(-1048)\cdot 42/5}{:}T\Lambda_Z{:}$,
\end{enumerate}
with $c+22/5=-1048/5\ne 0$ guaranteeing regularity of the rational
projector at $c=-214$.
\end{theorem}

\begin{proof}
Pope--Romans--Shen $1990$ NPB $339$ Eqs.~$(2.7$--$2.9)$ and
Bershadsky--Ooguri $1989$ CMP $126$ \S4 project each quasi-primary
onto the $\mathcal W_\infty[c]$-primary by subtracting
$\Lambda_Z$-descendants with scalar $-c/22$; at $c=-214$ this scalar
is $214/22=107/11$. At weight $4$, the projection is
$W_4=e_4+\tfrac{107}{11}\Lambda_Z$, equivalent to (i). At weight $5$,
three-leg quasi-primary uniqueness (Wang $1998$ CMP $195$ Prop.~$4.2$)
forces $e_5=W_5$ identically: the weight-$5$ $\mathcal W_\infty$-primary
has no $\Lambda_Z$-descendant with matching $L_1$-trivial projection.
At weight $6$, two $\Lambda_Z$-descendants appear ($\partial^2\Lambda_Z$
and ${:}T\Lambda_Z{:}$), each with projection coefficient read off
from the Pope--Romans--Shen $1990$ recursive primary formulas.
\end{proof}

\begin{proposition}[Three-path non-vanishing of $[e_k]$ at $c=-214$]
\label{prop:ek-nonvanish-vol3}
\ClaimStatusProvedHere
For $k=4,5,6$, $[e_k]\ne 0$ in $\ChirHoch^3(\mathbf H_{\Delta_5})_k$
at $c=-214$.
\end{proposition}

\begin{proof}
Three independent verification paths. (Path~A) The pairings
$\langle[\chi_k],[e_k]\rangle_{\Phi_k}\in
\{65193/10,5671/35,7704/7\}\cdot\mathrm{Vol}(E)(2\pi\mathrm i)^k$ of
Thms.~\ref{thm:chi4-e4-pairing-vol3}--\ref{thm:chi5-e5-pairing-vol3}
(and the analogous weight-$6$ computation) are nonzero rationals
times $\mathrm{Vol}(E)\ne 0$; non-vanishing of the pairing forces
non-vanishing of the class. (Path~B) Oberdieck--Pixton $2018$
reduced-DT polar-leading coefficients
$[p^{-1}q^{-1}y^k](-\Phi_{10}^{-1})$ reproduce the same rationals
modulo the Kapranov--Vasserot $\chi(\mathcal O_{\mathrm K3})=2$
normalisation, providing a curve-counting witness independent of
the Schiffmann--Vasserot CoHA path. (Path~C) Kontsevich $2003$ CY-$2$
formality scaling forces the leading behaviour to be linear in
$\mathrm{Vol}(E)$; the three paths agree at leading order.
\end{proof}

\subsection{Pattern for $e_k$ at general weight}
\label{subsec:ek-pattern}

\begin{conjecture}[Generic-weight $e_k$ shape]
\label{conj:ek-pattern-vol3}
\ClaimStatusConjectured
For $k\ge 4$, the chiral-Hochschild cocycle $e_k$ at $c=-214$ takes
the quasi-primary shape
\begin{equation}
\label{eq:ek-pattern-vol3}
e_k(z)={:}T\partial^{k-2}T{:}+\sum_{j=0}^{\lfloor(k-4)/2\rfloor}
\alpha_{k,j}\,{:}(\partial^{j+1}T)(\partial^{k-3-j}T){:}
+\beta_k\,\partial^kT+\hbar\,\mathrm{qt}(J^{(k)}),
\end{equation}
with leading tail $\beta_k=-c\,p_k(c)/q_k$ for $p_k$ a monic integer
polynomial of degree $\lfloor k/2\rfloor -1$ and $q_k$ a specific
integer, determined recursively by the $L_1$-quasi-primary condition.
For $k=4,5,6$ the explicit values are $\beta_4=-3c/20$,
$\beta_5=-c(c+2)/280$ (with a three-leg leading composite),
$\beta_6=-c(c-2)/42$.
\end{conjecture}

\emph{Primary.}
Zamolodchikov $1985$ TMF $65$;
Bershadsky--Ooguri $1989$ CMP $126$;
Pope--Romans--Shen $1990$ NPB $339$;
Bouwknegt--Schoutens $1993$ Phys.\ Rep.\ $223$;
Wang $1998$ CMP $195$;
Nakayama $2004$ IJMPA $19$;
Schiffmann--Vasserot $2013$ Publ.\ Math.\ IHES $118$;
Oberdieck $2018$ Geom.\ \& Topol.\ $22$;
Oberdieck--Pixton $2018$ Invent.\ Math.\ $213$;
Kontsevich $2003$ Lett.\ Math.\ Phys.\ $66$.
